% gerado por ferramentas/md2tex_enunciado.py a partir de questoes/q095.md (não editar à mão)
A mudança da resistência elétrica de alguns materiais em função da temperatura é utilizada para construir sensores que podem medir desde a temperatura ambiente até a de fornos. A figura representa um circuito elétrico usado para fazer a correspondência entre temperatura e tensão elétrica, na qual a resistência elétrica $R_s$ varia com a temperatura segundo a expressão $R_s = 20[1 + 0{,}002(T - 300)]$. A temperatura é medida em Kelvin.\par
\figq[0.37\linewidth]{q095_fig1.png}
A tensão elétrica, em volt, entre os pontos \textbf{A} e \textbf{B} da figura, quando o sensor estiver à temperatura de 800\,K é de\par
\begin{altlinha}
\alth{A}{5,0\,V.}%
\alth{B}{4,0\,V.}%
\alth{C}{3,0\,V.}%
\alth{D}{2,0\,V.}%
\alth{E}{1,0\,V.}
\end{altlinha}
